%% kapa: exam=1 round=27 year=2016 period=2 form=A subject=appraisal_law subject_ko=감정평가관계법규 seq=38 number=38
%% source: Q-Net 감정평가사 기출문제 (cst00309) artlSeq=5169402 file=extracted/1008/5169402/제27회 감정평가사 1차 문제/제27회 감정평가사 2교시문제(A형).pdf
%% conversion: kordoc 4.24.51 -> tools/gen1.py (auto); review=proofread-sonnet; flags=none
\begin{question}{38}
도시 및 주거환경정비법령상 정비구역 안에서 시장ㆍ군수의 허가를 받아야 하는 행위로 옳은 것만을 모두 고른 것은? (단, 재해복구 또는 재난수습에 필요한 응급조치를 위하여 하는 행위는 고려하지 않으며, 정비구역의 지정 및 고시 당시 이미 행위허가를 받았거나 받을 필요가 없는 행위는 제외함)
\begin{framedtext}
ㄱ. 가설공연장의 용도변경 \par
ㄴ. 죽목의 벌채 \par
ㄷ. 토지분할 \par
ㄹ. 이동이 용이하지 아니한 물건을 3주일 동안 쌓아놓는 행위 \par
\end{framedtext}
\begin{choices}
  \item ㄱ, ㄹ
  \item ㄷ, ㄹ
  \item ㄱ, ㄴ, ㄷ
  \item ㄱ, ㄷ, ㄹ
  \item ㄱ, ㄴ, ㄷ, ㄹ
\end{choices}
\end{question}
\answer{3}
